\documentclass[12pt,a4paper]{ctexart}
\usepackage[top=2.5cm,bottom=2.5cm,left=3cm,right=2.5cm]{geometry}
\usepackage{graphicx,float,booktabs,longtable,hyperref,xcolor}
\usepackage{fancyhdr,enumitem}
\graphicspath{{../screenshots/}}
\setCJKmainfont{SimSun}[BoldFont=SimHei]
\setmainfont{Times New Roman}
\pagestyle{fancy}
\fancyhead[L]{\small 尖兵项目2024C01066 — 验收审计综合材料}
\fancyhead[R]{\small 杭州数蛙科技有限公司}
\fancyfoot[C]{\thepage}
\renewcommand{\headrulewidth}{0.5pt}

\title{{\heiti\zihao{2} 超大规模智能物联网可信接入与数据治理关键技术研究 \\[6pt]}
       {\kaishu\large 合同编号: 2024C01066 | 课题五: 成果场景示范应用} \\[12pt]
       {\heiti\zihao{3} 验收审计综合材料}}
\author{杭州数蛙科技有限公司}
\date{\today}

\begin{document}
\maketitle\thispagestyle{empty}\newpage
\tableofcontents\newpage

% ===== Part 1 =====
\part{采购内容与项目关联性说明}

\section{概述}
杭州数蛙科技有限公司作为尖兵项目"超大规模智能物联网可信接入与数据治理关键技术研究"的参研单位，承担课题五"成果场景示范应用"的研发与交付任务。本文档系统阐述全部技术投入与项目的关联关系。

\section{技术投入总览}
\begin{table}[H]\centering
\caption{有合同采购项}\begin{tabular}{@{}clll@{}}
\toprule
序号 & 采购内容 & 供应商 & 合同文件 \\
\midrule
1 & 区块链设备(服务器+网关+安全模块) & 硬件供应商 & 硬件采购合同 \\
2 & 物联网电子签名管理系统V1.0 & 浙江爱签数字科技有限公司 & 软件采购合同 \\
\bottomrule
\end{tabular}\end{table}

\begin{table}[H]\centering
\caption{自主研发交付项}\begin{tabular}{@{}cll@{}}
\toprule
序号 & 研发项目 & 核心技术指标 \\
\midrule
1 & 时序数据采集管理系统 & 10+协议/物模型TSL/边缘流式计算/驾驶舱可视化 \\
2 & AIP运维管理平台 & 8视图/OWL本体836三元组/SPARQL/ISA-95对齐 \\
3 & 技术集成与联调服务 & 11个示范场景部署/全系统联调/87项测试 \\
\bottomrule
\end{tabular}\end{table}

\section{需求维度关联}
\begin{table}[H]\centering
\caption{投入内容与项目需求对标}
\begin{tabular}{@{}p{3cm}p{5cm}p{5cm}@{}}
\toprule
投入内容 & 对标的合同需求 & 不投入的后果 \\
\midrule
区块链设备(采购合同) & 联盟链网络，PBFT共识，可信接入基础设施 & 无法构建区块链网络 \\
电子签名系统(爱签合同) & IoT数据电子签名及区块链上链存证 & 无合规签名机制 \\
时序采集系统(自研) & 10+协议适配/物模型/边缘计算/驾驶舱 & 核心采集能力缺失 \\
AIP运维平台(自研) & 8视图运维/OWL本体/SPARQL & 运维管理能力缺失 \\
\bottomrule
\end{tabular}\end{table}

\section{技术实现指标关联}
\begin{table}[H]\centering
\caption{投入对技术指标的支撑}
\begin{tabular}{@{}p{3.5cm}p{5cm}p{5cm}@{}}
\toprule
投入内容 & 支撑的技术指标 & 技术实现方式 \\
\midrule
区块链设备 & 链上TPS$\ge$1,000; 拜占庭容错 & PBFT/Hyperledger Fabric 2.5/HSM \\
电子签名系统 & SM2国密/区块链存证/不可篡改 & 爱签V1.0 API/SM2+SM3+SM4 \\
时序采集系统 & 10+协议/毫秒级采集/10万+并发 & 插件注册制/滑动窗口/TDengine超级表 \\
AIP运维平台 & 836三元组/ISA-95/SPARQL & FastAPI+OWL+W3C标准 \\
\bottomrule
\end{tabular}\end{table}

\section{合同技术规范逐条对标}
\begin{table}[H]\centering
\caption{合同条款与科技报告对标}
\begin{tabular}{@{}p{4cm}p{4cm}p{3cm}@{}}
\toprule
采购合同技术规范 & 科技报告对应条款 & 结果 \\
\midrule
区块链设备支持PBFT,$\ge$8节点 & Y.2.1 联盟链PBFT共识 & 一致 \\
电子签名支持SM2及区块链存证 & Y.4.1 IoT上链国密认证 & 一致 \\
TDengine 10万+测点毫秒级写入 & X.3.2 时序库性能要求 & 一致 \\
DG-IoT自研MQTT引擎(百万级并发+Sparkplug B) & X.2.3 大规模设备接入 & 一致 \\
硬件兼容国产CPU+麒麟V10 & 信创要求 & 一致 \\
\bottomrule
\end{tabular}\end{table}

\newpage
% ===== Part 2 =====
\part{自研软件审计证据链}

\section{代码开发持续记录（Git Commit Log摘要）}
\begin{table}[H]\centering\caption{核心版本节点}\begin{tabular}{@{}llllll@{}}
\toprule
日期 & 作者 & 模块 & 说明 & 文件数 & 代码行 \\
\midrule
2024-03-15 & dev01 & ontology.py & OWL 5层本体初始化 & 3 & +1,200 \\
2024-05-10 & dev01 & graphrag.py & GraphRAG搜索引擎 & 1 & +1,500 \\
2024-08-10 & dev01 & safety\_rules.py & 26条SWRL规则 & 2 & +1,100 \\
2024-10-20 & dev03 & BlockchainView & 区块链8Tab平台 & 5 & +1,800 \\
2024-12-15 & dev01 & sparql & W3C SPARQL端点 & 1 & +350 \\
2025-02-10 & dev03 & AIP 8视图 & Explorer/Ontology等 & 8 & +2,400 \\
2025-04-20 & dev02 & opcua.py & OPC UA适配器 & 1 & +480 \\
2025-06-15 & dev03 & test\_graphrag & 87项自动化测试 & 4 & +1,200 \\
\bottomrule
\end{tabular}\end{table}
\vspace{6pt}
提交统计: 20个版本节点 | 4名核心开发 | 75+源文件 | 约20,000行代码 | 持续21个月

\section{项目组员工时台账}
\begin{table}[H]\centering
\caption{工时汇总（费率依据CSBMK-202510）}
\begin{tabular}{@{}llrrrr@{}}
\toprule
角色 & 姓名 & 人天数 & 人天单价(元) & 金额(万元) \\
\midrule
高级全栈工程师 & developer01 & 95 & 1,200 & 11.4 \\
协议工程师 & developer02 & 85 & 1,100 & 9.35 \\
前端工程师 & developer03 & 90 & 1,100 & 9.9 \\
全栈工程师 & developer04 & 80 & 1,100 & 8.8 \\
项目经理 & project\_manager & 25 & 1,200 & 3.0 \\
测试工程师 & qa\_engineer & 40 & 1,000 & 4.0 \\
\midrule
合计 & & 415 & & 46.25 \\
\bottomrule
\end{tabular}\end{table}

\section{功能点量化分析}
依据DB3303/T 059-2023《政务信息化项目软件开发费用测算规范》：
\begin{table}[H]\centering\caption{功能点统计}\begin{tabular}{@{}lrrr@{}}
\toprule
功能点类型 & 数量 & 复杂度 & 调整后FP \\
\midrule
内部逻辑文件(ILF) & 10 & 中-高 & 110 \\
外部接口文件(ELF) & 4 & 高 & 60 \\
外部输入(EI) & 6 & 中 & 36 \\
外部输出(EO) & 8 & 中 & 48 \\
外部查询(EQ) & 5 & 中 & 25 \\
\midrule
合计 & 35 & & 279 \\
\bottomrule
\end{tabular}\end{table}
279功能点 $\times$ 1,500元/FP $\approx$ 41.85万元。加技术集成联调(5万)及部署验收(约8万)，合计约55万元。

\section{行业基准数据}
依据《2025年中国软件行业基准数据》（CSBMK-202510）：
\begin{table}[H]\centering\caption{杭州软件开发基准}\begin{tabular}{@{}lll@{}}
\toprule
基准指标 & 杭州数据 & 来源 \\
\midrule
人月费率(P50) & 28,813元/人月 & CSBMK-202510 \\
人天单价 & 1,325元/人天 & 计算值(21.75天/月) \\
生产率 & 8-12 FP/人月 & CSBMK-202510 \\
\bottomrule
\end{tabular}\end{table}
本项目采用人天单价1,000-1,200元，均低于行业基准，费用估算保守合理。

\section{政策依据}
\begin{enumerate}[leftmargin=*]
\item 浙财科教〔2025〕38号《浙江省科技发展专项资金管理办法》
\item 浙注协〔2025〕185号《浙江省科技计划项目专项审计报告编制指引》
\item DB3303/T 059-2023《政务信息化项目软件开发费用测算规范》
\item CSBMK-202510《2025年中国软件行业基准数据》
\end{enumerate}

\newpage
% ===== Part 3 =====
\part{系统运行截图留档}

\begin{figure}[H]\centering
\includegraphics[width=0.95\textwidth]{dashboard.png}
\caption{仪表盘首页 — IoT数据采集 + 区块链KPI + 尖兵项目总览}
\end{figure}
\begin{figure}[H]\centering
\includegraphics[width=0.95\textwidth]{blockchain.png}
\caption{区块链可信管理平台 — 8大功能面板（含电子签名存证）}
\end{figure}
\begin{figure}[H]\centering
\includegraphics[width=0.95\textwidth]{aip.png}
\caption{AIP运维管理平台 — 8视图（Explorer/Ontology/SPARQL/Pipeline/Console/Tree/Graph/Validate）}
\end{figure}

\newpage
\part*{结语}
\addcontentsline{toc}{part}{结语}

数蛙在尖兵项目中的全部技术投入均与课题五"成果场景示范应用"存在直接且不可替代的关联关系。所有采购合同有据可查，自研软件有Git提交日志、工时台账、功能点分析和行业基准数据四重交叉验证，技术系统可现场演示。上述投入是保障课题五按期完成研发交付的必要前提。

\vspace{12pt}
\noindent\textbf{编制单位}：杭州数蛙科技有限公司 \hfill \textbf{编制日期}：\today

\end{document}
